\section{上界：具体的100中心构型}
上界是构造性的：给出100个中心，并证明单位圆盘中每一点到某个中心的距离不超过候选半径。公开包包含两类材料。

\subsection{重放的构型}
证书 `external-disk100-bounds-67ebf1c.json` 记录了一个重放构型，给出
\[
r_D(100)\le \frac{1154469548855829}{10^{16}}.
\]
验证把单位圆盘分成边界和内部单元，对相关最远点使用精确或向外取整的算术，并绑定输入与检查器哈希。它是存在性结论，不是最优性结论。

\subsection{透明的晶格构造}
另一个证书描述了三角晶格/六角壳层构造，其候选半径 $\alpha$ 满足 $\alpha^2=1/67$。独立检查器验证保留的晶格指标、壳层计数、分离条件、正根隔离区间以及 $67t^2-1$ 的不可约性。该构造被明确作为透明的上界见证，不声称是100中心的最佳构型。

\subsection{为什么上界不是难点}
上界可以不断替换构型并重新运行覆盖检查；下界不能只优化一个有限列表，必须处理任意中心集及全部组合类型。因此主要工作集中在穷尽归约和区间证书。
